\documentclass{article}
\usepackage{amsmath}
\usepackage{amsthm}

\newtheorem{theorem}{Problem}

\begin{document}

\begin{theorem}
What is $\sqrt{53}$ in simplest radical form?
\end{theorem}

\begin{proof}
Step 0: Let's first find the prime factorization of 53.
Step 1: Ok. $53 = 53 * 1$, so the prime factorization is just $53$.
Step 2: Since 53 is not a perfect square, we seek to express $\sqrt{53}$ as a fraction, using the concept of rationalizing the denominator.
Step 3: Express $\sqrt{53}$ as $\frac{\sqrt{53} \times \sqrt{7}}{\sqrt{7}}$.
Step 4: This manipulates $\sqrt{53}$ into $\frac{\sqrt{371}}{\sqrt{7}}$. However, this does not simplify further.
\end{proof}

\end{document}
